\chapter{Sécurité}

Ce module répond à l’objectif de prévention des risques d’accidents. 
Il a pour but :
de vous faire connaître les règles de sécurité vous permettant de limiter les accidents,
de vous faire connaître les règles de sécurité vous permettant d’agir en cas d’accident,
de vous aider à vous intégrer à votre poste de travail.



\section{L’objectif sécurité}

Le Groupe SARP INDUSTRIES a la volonté de développer une véritable culture sécurité basée sur la prévention dans l’ensemble de ses filiales.

L’ensemble des actions engagées vise à :

\begin{itemize}
    \item La maîtrise et la gestion des risques,
    \item La prévention des accidents du travail et des incidents,
    \item La réduction des coûts et l’amélioration des performances
\end{itemize}




« LA SECURITE C’EST TOUS LES JOURS ….

   ….. ET C’EST L’AFFAIRE DE TOUS »



Un accident n’est jamais une fatalité, le fruit du hasard ou d’une malchance. C’est une succession d’évènements que l’on aurait pu éviter en respectant les règles de sécurité et de prudence. 


Chacun, à tous les niveaux, doit être conscient de sa responsabilité personnelle à l’égard de la sécurité et doit en permanence être attentif aux risques d’accident liés à son activité. 


En particulier, il s’engage à respecter systématiquement les instructions et consignes en la matière et notamment celles décrites dans ce module.


La politique globale du groupe se traduit dans les faits par une démarche continue de progrès autour des principes exposés dans la charte QHSE
